Este capítulo documenta el artefacto evaluado. Las primeras secciones describen la arquitectura de JobAI, el proceso de ingesta y el pipeline de consulta en su configuración de referencia, que es la que se evalúa en el Capítulo~\ref{cap:resultados}. Las secciones finales corresponden a las dos iteraciones de diseño realizadas en este trabajo. La primera diagnosticó y migró el almacenamiento vectorial y corrigió los defectos detectados en la revisión del código. La segunda incorporó a la recuperación filtros estructurados, búsqueda híbrida, reordenamiento y un modelo de \textit{embeddings} multilingüe.

\section{Visión general de la arquitectura}
\label{sec:art-arquitectura}

JobAI es una plataforma de agregación de ofertas de empleo remoto cuyo componente central es un copiloto conversacional basado en RAG. El sistema se organiza en tres capas (Figura~\ref{fig:art-arquitectura}):
\begin{itemize}
    \item un cliente web en Next.js, responsable de la autenticación, la navegación del catálogo y el panel conversacional;
    \item un servidor en NestJS, que orquesta la ingesta, la recuperación y la generación;
    \item una capa de datos compuesta por PostgreSQL, como almacén relacional, y un almacén vectorial que, tras la iteración descrita en la Sección~\ref{sec:art-vectorial}, puede ser pgvector o Qdrant.
\end{itemize}
Los modelos se consumen como servicios. Un servidor propio, ejecutado localmente sobre una GPU, sirve el modelo de \textit{embeddings} y el modelo de reordenamiento; los modelos de lenguaje se consumen a través de un proveedor externo, Groq.

Este capítulo describe la \emph{configuración de referencia} del artefacto, definida en el repositorio junto con el código. El despliegue público de la plataforma conserva todavía la configuración anterior a la segunda iteración, con un proveedor externo de \textit{embeddings} y sin reordenamiento, y se actualizará a la configuración de referencia.

\begin{figure}[htbp!]
    \centering
    \caption{Arquitectura de JobAI}
    \label{fig:art-arquitectura}
    \begin{tikzpicture}[
        box/.style={draw, rounded corners, align=center, minimum height=1cm, text width=3.2cm, font=\small},
        store/.style={draw, cylinder, shape border rotate=90, aspect=0.15, align=center, minimum height=1.2cm, text width=2.4cm, font=\small},
        arr/.style={->, >=stealth}
    ]
        \node[box, text width=2.2cm] (fe) at (0,0) {Cliente web\\(Next.js)};
        \node[box, text width=2.2cm] (be) at (5.4,0) {Servidor\\(NestJS)};
        \node[box] (models) at (10.2,2.1) {Modelos de lenguaje\\(Groq)};
        \node[box] (local) at (10.2,0) {\textit{Embeddings} y reordenamiento (servidor local, GPU)};
        \node[box] (sources) at (10.2,-2.1) {Fuentes de ofertas\\(Remotive, Jobicy, Findwork)};
        \node[store] (pg) at (3.8,-3.2) {PostgreSQL\\(+ pgvector)};
        \node[store] (qd) at (7.0,-3.2) {Qdrant};
        \draw[arr] (fe) -- node[above, font=\scriptsize] {REST / NDJSON} (be);
        \draw[arr] (be) -- (models);
        \draw[arr] (be) -- (local);
        \draw[arr] (sources) -- (be);
        \draw[arr] (be) -- (pg);
        \draw[arr, dashed] (be) -- (qd);
    \end{tikzpicture}
\end{figure}

El servidor se organiza en módulos con responsabilidades separadas (Tabla~\ref{tab:art-modulos}). La consulta conversacional se expone mediante un \textit{endpoint} que transmite la respuesta de forma incremental como una secuencia de objetos JSON delimitados por saltos de línea (NDJSON). Los eventos son de cuatro tipos: inicio, \textit{token}, fin con las fuentes citadas y error.

\begin{table}[htbp!]
    \centering
    \caption{Módulos del servidor de JobAI}
    \label{tab:art-modulos}
    \begin{tabular}{|p{3.2cm}|p{11cm}|}
        \hline
        \textbf{Módulo} & \textbf{Responsabilidad} \\
        \hline
        Ingesta & Obtención periódica de ofertas, deduplicación, análisis estructurado, normalización de salarios y fragmentación \\
        \hline
        \textit{Embeddings} & Generación de vectores de consultas y fragmentos, con caché por modelo \\
        \hline
        Consulta & Clasificación de la consulta, extracción de restricciones y orquestación de las rutas de recuperación, agregación e híbrida \\
        \hline
        Recuperación & Búsqueda densa y léxica bajo el mismo filtro, fusión de rankings y reordenamiento \\
        \hline
        RAG & Construcción del contexto y del prompt \\
        \hline
        LLM & Acceso a los modelos generativos, con reintentos y rotación de claves \\
        \hline
        Almacenamiento & Repositorios relacionales y almacenes vectoriales (pgvector y Qdrant) \\
        \hline
        Currículo & Carga, extracción de texto y análisis estructurado del currículo del usuario \\
        \hline
        Autenticación y usuario & Tokens JWT con rotación, perfil y ofertas guardadas \\
        \hline
    \end{tabular}
\end{table}

\section{Ingesta de datos}
\label{sec:art-ingesta}

Un proceso programado consulta las tres fuentes públicas cada seis horas. Puede desactivarse mediante configuración, como se hizo para congelar el corpus de evaluación. Cada ciclo sigue cuatro pasos:
\begin{enumerate}
    \item \textbf{Deduplicación.} Cada oferta se identifica mediante el resumen SHA-256 de su título, empresa, descripción y ubicación. Las ofertas ya almacenadas y con \textit{embeddings} se omiten.
    \item \textbf{Análisis estructurado.} Las ofertas nuevas pasan por un modelo de lenguaje que extrae en formato JSON los siguientes campos: resumen, salario, responsabilidades, requisitos, beneficios, habilidades, tipo de contrato y modalidad. Ese modelo es Llama 3.1 8B local en desarrollo y \texttt{gpt-oss-120b} vía Groq en producción. El prompt instruye explícitamente al modelo a devolver \texttt{null} cuando un campo no aparece en el texto, en lugar de inferirlo. Los salarios se normalizan a dólares anuales; los valores fuera del intervalo de 10.000 a 1.000.000 dólares, o con un mínimo mayor que el máximo, se registran como no informados, sin intentar adivinar la moneda o el periodo.
    \item \textbf{Fragmentación.} Si el análisis estructurado produce al menos un campo, la oferta se divide en fragmentos \emph{tipados}:
    \begin{itemize}
        \item un fragmento de identidad (título, empresa, ubicación, modalidad y resumen);
        \item uno de requisitos;
        \item uno de responsabilidades y habilidades;
        \item uno de beneficios y salario.
    \end{itemize}
    Si el análisis falla, se aplica una ventana deslizante de 800 caracteres con solapamiento de 200, que respeta los límites de párrafo, viñeta y oración, con un máximo de cuatro ventanas por oferta.
    \item \textbf{\textit{Embeddings} y almacenamiento.} Cada fragmento se convierte en un vector de 768 dimensiones con el modelo multilingual-e5-base \citep{Wang2024mE5}, con el prefijo \texttt{passage:} que el modelo espera para los documentos, y se almacena junto con su texto, su tipo y el modelo que lo produjo. El servidor rechaza vectores de un modelo distinto al configurado, de modo que una consulta nunca se compara con fragmentos codificados por otro modelo.
\end{enumerate}
Al terminar cada ciclo se reconstruye el índice léxico que usa la búsqueda híbrida (Sección~\ref{sec:art-recuperacion}).

La fragmentación tipada responde a un hallazgo de una iteración previa (Sección~\ref{sec:met-ciclos-previos}). Un fragmento que mezcla requisitos, beneficios y descripción compite peor en la recuperación que un fragmento que contiene solo la información por la que se pregunta.

\section{Pipeline de consulta}
\label{sec:art-consulta}

La Figura~\ref{fig:art-pipeline} resume el recorrido de una consulta.

\begin{figure}[htbp!]
    \centering
    \caption{Pipeline de consulta de JobAI}
    \label{fig:art-pipeline}
    \begin{tikzpicture}[
        stage/.style={draw, rounded corners, align=center, text width=4.2cm, minimum height=0.9cm, font=\small},
        half/.style={stage, text width=3.2cm},
        decision/.style={draw, diamond, aspect=2.2, align=center, inner sep=1pt, font=\small},
        arr/.style={->, >=stealth}
    ]
        \node[stage] (q) at (0,0) {Consulta del usuario};
        \node[decision] (cls) at (0,-1.6) {Clasificación};
        \node[stage] (rag) at (-3.6,-3.2) {Ruta de recuperación\\(con restricciones como filtro)};
        \node[stage] (agg) at (3.6,-3.2) {Ruta de agregación\\(plantilla SQL)};
        \node[half] (dense) at (-5.5,-4.8) {Búsqueda densa: 50 fragmentos $\rightarrow$ ofertas};
        \node[half] (bm25) at (-1.7,-4.8) {BM25:\\50 ofertas};
        \node[stage] (rrf) at (-3.6,-6.4) {Fusión RRF\\(20 candidatas)};
        \node[stage] (rr) at (-3.6,-7.8) {Reordenamiento\\(5 ofertas)};
        \node[stage] (ctx) at (-3.6,-9.2) {Contexto por oferta + campos estructurados};
        \node[stage] (sum) at (3.6,-4.8) {Resumen de la agregación};
        \node[stage] (llm) at (0,-10.8) {Prompt + generación (LLM)};
        \node[stage] (out) at (0,-12.2) {Respuesta incremental con fuentes};
        \draw[arr] (q) -- (cls);
        \draw[arr] (cls) -| node[above, pos=0.25, font=\scriptsize] {recuperación} (rag);
        \draw[arr] (cls) -| node[above, pos=0.25, font=\scriptsize] {agregación} (agg);
        \draw[arr] (rag) -- (dense);
        \draw[arr] (rag) -- (bm25);
        \draw[arr] (dense) -- (rrf);
        \draw[arr] (bm25) -- (rrf);
        \draw[arr] (rrf) -- (rr);
        \draw[arr] (rr) -- (ctx);
        \draw[arr] (agg) -- (sum);
        \draw[arr] (ctx) |- (llm);
        \draw[arr] (sum) |- (llm);
        \draw[arr] (llm) -- (out);
    \end{tikzpicture}
\end{figure}

\subsection{Clasificación de la consulta}
\label{sec:art-clasificacion}

Un conjunto de patrones \textit{regex} detecta señales léxicas de tres tipos:
\begin{itemize}
    \item de recuperación (\textit{find}, \textit{looking for}, \textit{jobs with});
    \item de agregación (\textit{how many}, \textit{average}, \textit{salary}, \textit{list all});
    \item de restricción, en inglés y en español: montos de salario, modalidad (\textit{remote}, \textit{remoto}, \textit{presencial}) y tipo de contrato (\textit{contract}, \textit{tiempo completo}, \textit{freelance}, \textit{prácticas}).
\end{itemize}
Las consultas con señales de recuperación y sin señales de agregación ni de restricción van directamente a la ruta de recuperación, sin invocar ningún modelo de lenguaje. En los demás casos, un clasificador basado en LLM, con temperatura 0, decide entre recuperación, agregación o una ruta híbrida. Según el caso, el clasificador selecciona además una de 24 plantillas SQL parametrizadas, como conteos, estadísticas salariales o listados por habilidad, o extrae las restricciones de la consulta: salario mínimo o máximo en dólares anuales, modalidad remota y tipo de contrato. El modelo nunca genera SQL libre, y las consultas de agregación se ejecutan con una conexión de solo lectura.

Sobre la respuesta del clasificador se aplican tres reglas de enrutamiento deterministas:
\begin{itemize}
    \item una consulta con restricciones extraídas se atiende por la ruta de recuperación, con las restricciones como filtro;
    \item las plantillas que responden con estadísticas, valores distintos u ofertas recientes se atienden siempre por la ruta de agregación, aunque el clasificador haya etiquetado la consulta como recuperación;
    \item una consulta de agregación sin una plantilla válida recurre a la ruta de recuperación.
\end{itemize}

El título del cargo y las ubicaciones distintas de \textit{remoto} no se convierten en filtros: exigir su coincidencia literal reproduciría la limitación léxica que el pipeline busca superar, así que se dejan a la recuperación semántica. Como las fuentes suelen etiquetar el trabajo independiente como contrato, una consulta por trabajo \textit{freelance} acepta ambos tipos. Los filtros que el usuario fija explícitamente en el panel de búsqueda tienen prioridad sobre los extraídos del texto.

\subsection{Recuperación}
\label{sec:art-recuperacion}

La recuperación combina dos búsquedas de primera etapa, una fusión de rankings y un reordenamiento:
\begin{enumerate}
    \item \textbf{Búsqueda densa.} La consulta se convierte en un vector con el prefijo \texttt{query:}. Si el usuario cargó un currículo, se antepone a la consulta un perfil breve con sus tres cargos más recientes y sus diez habilidades principales. Se buscan los 50 fragmentos más similares por coseno que cumplen el filtro; se agrupan por oferta y cada oferta toma la similitud de su mejor fragmento.
    \item \textbf{Búsqueda léxica.} En paralelo, BM25 \citep{Robertson2009}, con $k_1 = 1{,}2$ y $b = 0{,}75$, busca las 50 ofertas que mejor coinciden con el texto de la consulta entre las que cumplen el mismo filtro. El índice se mantiene en memoria sobre el título, las habilidades y la descripción de cada oferta, y es la misma implementación que sirve de línea base en la evaluación.
    \item \textbf{Fusión.} Las dos listas se combinan mediante fusión por rangos recíprocos (RRF) \citep{Cormack2009}: cada oferta suma $1/(k + r)$ por cada lista en la que aparece en la posición $r$, con $k = 60$. Como la fusión usa solo las posiciones, no es necesario calibrar entre sí la similitud coseno y el puntaje de BM25.
    \item \textbf{Reordenamiento.} Las 20 primeras ofertas de la lista fusionada pasan por un modelo de reordenamiento de tipo \textit{cross-encoder}, Qwen3-Reranker-0.6B \citep{Zhang2025Qwen3}, que lee juntas la consulta y la oferta \citep{Nogueira2019} y estima la probabilidad de que la oferta sea relevante. El modelo recibe el encabezado de la oferta (título, empresa, ubicación, tipo de contrato y salario) y sus fragmentos, hasta 2.000 caracteres. Se conservan las cinco ofertas con mayor probabilidad.
\end{enumerate}

El filtro se aplica dentro de cada motor, durante la búsqueda y no sobre sus resultados: Qdrant lo evalúa sobre la carga útil de los puntos, pgvector lo incorpora a la consulta SQL, con barridos iterativos cuando usa el índice, y BM25 descarta las ofertas que no lo cumplen antes de cortar en 50. Una oferta sin salario informado nunca cumple una condición de salario. Si ninguna oferta cumple el filtro, la búsqueda se repite sin él y el prompt indica al modelo que debe advertir al usuario que ninguna oferta satisface todas sus condiciones.

Dos mecanismos completan la etapa. El primero es una compuerta de \emph{sin resultados}: si incluso la mejor oferta reordenada tiene una probabilidad menor que 0,134, el sistema responde que no encontró ofertas relevantes en lugar de presentar resultados que no lo son. El umbral se calibró sobre un conjunto de consultas de desarrollo (Sección~\ref{sec:art-iter-recuperacion}). El segundo es la degradación ante fallos: si el servidor de reordenamiento no responde, se conserva el orden de la fusión.

\subsection{Construcción del contexto}

Para cada una de las cinco ofertas se arma un bloque de contexto con un encabezado (título, empresa, ubicación, tipo de contrato, salario y similitud) y el texto de sus fragmentos, limitado a 2.400 caracteres.

Antes de cerrar el contexto, una etapa de \emph{detección de campos} identifica si la consulta pregunta por salario, beneficios, requisitos, habilidades o responsabilidades. Si el campo solicitado no aparece en los fragmentos recuperados, se añade desde los datos estructurados de la oferta. Esta etapa reduce el riesgo de que el modelo complete con información plausible pero inexistente un dato que sí está en la base de datos y simplemente no fue recuperado.

\subsection{Generación}

El prompt de sistema instruye al modelo a responder \emph{únicamente} a partir de las ofertas proporcionadas, a indicar explícitamente cuando estas no contienen la información solicitada, a no inventar detalles y a responder en el idioma de la pregunta aunque las ofertas estén en otro. El historial reciente de la conversación, hasta 12 mensajes, se incluye para resolver preguntas de seguimiento; la conversación solo afecta a la generación y no a la recuperación, que usa la consulta actual. La generación usa el modelo \texttt{qwen/qwen3.8-27b} servido por Groq, con temperatura 0,2 en producción y respuesta incremental. En la ruta híbrida, los resultados de la recuperación y de la agregación se ejecutan en paralelo y se combinan en un único prompt; si una de las dos rutas falla, se responde con la otra.

\section{Iteración de diseño: almacenamiento vectorial}
\label{sec:art-vectorial}

\subsection{Diagnóstico del estado inicial}

En su versión inicial, JobAI almacenaba los vectores en una columna \texttt{vector(768)} de PostgreSQL, mediante la extensión pgvector \citep{Pgvector}. El esquema incluía un índice IVFFlat creado en la migración inicial, es decir, sobre una tabla todavía vacía y con los parámetros por defecto. La consulta de similitud, sin embargo, ordenaba los resultados por la expresión derivada $1 - d(\mathbf{q}, \mathbf{v})$ y filtraba por umbral sobre esa misma expresión. El planificador de PostgreSQL solo puede servir una búsqueda de vecinos cercanos desde el índice cuando la consulta ordena directamente por el operador de distancia (\texttt{ORDER BY embedding <=> q LIMIT k}).

La Tabla~\ref{tab:art-explain} resume la inspección con \texttt{EXPLAIN ANALYZE} sobre el corpus de evaluación (4.013 fragmentos), antes y después de la migración del índice.

\begin{table}[htbp!]
    \centering
    \caption{Plan de ejecución de la búsqueda de 15 vecinos en pgvector sobre el corpus de evaluación}
    \label{tab:art-explain}
    \begin{tabular}{|p{5.2cm}|p{2.2cm}|p{4.6cm}|r|}
        \hline
        \textbf{Consulta} & \textbf{Índice} & \textbf{Plan elegido} & \textbf{Tiempo} \\
        \hline
        Original (orden por $1-d$) & IVFFlat & Barrido secuencial + ordenamiento & 25,7 ms \\
        \hline
        Reescrita (orden por $d$) & IVFFlat & Recorrido del índice & 2,8 ms \\
        \hline
        Reescrita (orden por $d$) & HNSW & Barrido secuencial + ordenamiento & 24,8 ms \\
        \hline
        Reescrita, índice forzado & HNSW & Recorrido del índice & 2,0 ms \\
        \hline
    \end{tabular}
\end{table}

La inspección confirma tres hechos.
\begin{itemize}
    \item La consulta original nunca usaba el índice: recorría y ordenaba la tabla completa.
    \item El índice IVFFlat creado sobre la tabla vacía no perdía \textit{recall} (1,000 frente a la búsqueda exacta en 100 consultas), probablemente porque todos los vectores quedaron en una sola lista; en la práctica se comportaba como un recorrido completo.
    \item Con el índice HNSW y la consulta ya corregida, el planificador de PostgreSQL sigue prefiriendo el barrido secuencial a esta escala, aunque el recorrido del índice es 12 veces más rápido. El modelo de costos de pgvector sobrestima el costo de arranque del índice en tablas de pocos miles de filas. Por eso la búsqueda admite un modo que fuerza el índice, y la evaluación mide ambos comportamientos.
\end{itemize}

La misma inspección reveló que el umbral de similitud de 0,5 no filtra nada con el modelo de entonces, e5-base-v2. En 20 consultas de prueba, la similitud coseno mínima entre la consulta y \emph{cualquier} fragmento del corpus fue 0,66, y la mediana 0,75. Las similitudes de e5-base-v2 se concentran en valores altos, así que todos los fragmentos superan el umbral. La segunda iteración reemplazó este umbral por la compuerta de sin resultados descrita en la Sección~\ref{sec:art-recuperacion}.

\subsection{Diseño de la migración}

La migración se diseñó para cumplir dos requisitos: que la plataforma pueda servirse desde Qdrant \citep{QdrantDocs}, y que ambos motores puedan compararse en igualdad de condiciones. Las decisiones principales fueron cinco.

\begin{description}
    \item[Abstracción del almacén vectorial.] Se definió una interfaz única con estas operaciones:
    \begin{itemize}
        \item escritura de los fragmentos de una oferta;
        \item búsqueda con número de resultados, umbral, filtros estructurados, modo y parámetro $ef_{search}$;
        \item búsqueda restringida a un conjunto de ofertas;
        \item verificación de existencia y borrado.
    \end{itemize}
    Los módulos de ingesta, RAG y gestión de ofertas dependen de la interfaz y no de una implementación concreta. Una variable de configuración selecciona el motor que sirve a la aplicación.
    \item[Corrección de la consulta en pgvector.] La búsqueda se reescribió para ordenar por el operador de distancia, y el umbral se aplica después de obtener los $k$ vecinos, lo cual es equivalente porque la similitud decrece monótonamente con la distancia. El índice IVFFlat se reemplazó por un índice HNSW \citep{Malkov2020} con $m = 16$ y $ef_{construction} = 128$. La búsqueda admite tres modos:
    \begin{itemize}
        \item \emph{automático}: decide el planificador, que en tablas pequeñas puede preferir el barrido secuencial;
        \item \emph{aproximado}: fuerza el uso del índice;
        \item \emph{exacto}: fuerza la fuerza bruta, que es la referencia para medir el \textit{recall}.
    \end{itemize}
    Para búsquedas con filtros se activan los barridos iterativos del índice, mediante el parámetro \texttt{hnsw.\allowbreak iterative\_scan}. Sin ellos, un filtro selectivo puede devolver menos de $k$ resultados.
    \item[Colección en Qdrant.] La colección usa distancia coseno y los mismos parámetros HNSW que pgvector. Cada punto almacena como carga útil (\textit{payload}) el texto y el tipo del fragmento y los atributos de la oferta usados en los filtros: salario, tipo de contrato, modalidad y fuente, con índices de \textit{payload} sobre los campos filtrables. Los identificadores de fragmento de PostgreSQL no son UUID, por lo que el identificador de cada punto se deriva de forma determinista mediante un UUID versión 5. Se redujo además el umbral a partir del cual Qdrant construye el índice HNSW de un segmento: con el valor por defecto, un corpus de unos pocos miles de vectores se resolvería por fuerza bruta y la comparación no mediría el índice.
    \item[PostgreSQL como fuente de verdad.] El texto de los fragmentos y su relación con las ofertas permanecen en PostgreSQL, con borrado en cascada e integridad referencial. Al escribir, el almacén de Qdrant primero confirma la transacción en PostgreSQL y luego reemplaza los puntos de la oferta. La verificación de si una oferta ya tiene \textit{embeddings} se hace contra Qdrant, de modo que una escritura fallida en Qdrant se repara en el siguiente ciclo de ingesta.
    \item[Carga y verificación.] Un comando de línea copia en lotes todos los vectores existentes de PostgreSQL a Qdrant. La operación es idempotente, porque los identificadores son deterministas. Al terminar, el comando espera a que Qdrant indexe la colección y verifica que coincidan tres conteos: filas en PostgreSQL, puntos en Qdrant y vectores indexados.
\end{description}

Sobre el corpus de evaluación, la búsqueda exacta en ambos motores coincidió en el 99,85\,\% de los 10 primeros resultados de 200 consultas. Las diferencias corresponden a empates entre fragmentos de texto idéntico, que cada motor ordena de forma distinta. Esto confirma que los dos almacenes contienen los mismos datos y usan la misma métrica. La comparación de desempeño entre ambos motores se reporta en el Capítulo~\ref{cap:resultados}.

\section{Defectos detectados en la revisión previa a la evaluación}
\label{sec:art-defectos}

Antes de evaluar el artefacto se revisó su código frente a su documentación. La revisión encontró cinco defectos, que se corrigieron dentro de esta iteración porque habrían invalidado parte de la evaluación:

\begin{enumerate}
    \item \textbf{La personalización por currículo reemplazaba la consulta.} La documentación del artefacto indicaba que el perfil se antepone a la consulta del usuario. La implementación, en cambio, construía la búsqueda \emph{solo} con el perfil, de modo que un usuario con currículo recibía las mismas ofertas fuera cual fuera su pregunta. Se corrigió para anteponer el perfil.
    \item \textbf{La ruta de agregación fallaba con las consultas que solo contenían señales de agregación.} Los patrones \textit{regex} clasificaban la consulta como agregación sin seleccionar ninguna plantilla SQL, y la ejecución fallaba al no encontrar la plantilla. Las pruebas de integración no lo detectaban porque simulaban el repositorio SQL. Se corrigió delegando la selección de la plantilla al clasificador LLM, con recuperación como alternativa; las pruebas ahora verifican la plantilla seleccionada.
    \item \textbf{Las consultas de similitud no usaban el índice vectorial.} Se describe en la Sección~\ref{sec:art-vectorial}.
    \item \textbf{Las plantillas de agregación sobre ofertas remotas ignoraban el atributo de modalidad.} Las plantillas que cuentan y listan ofertas remotas buscaban la palabra \textit{remote} en el texto de la ubicación. No usaban el atributo booleano de modalidad que la ingesta obtiene de la fuente o del análisis estructurado. Sobre el corpus de evaluación, la plantilla original reconocía como remotas 39 ofertas, frente a más de 570 marcadas como tales. Se corrigió para usar el atributo, conservando el texto de la ubicación como respaldo.
    \item \textbf{Otros defectos menores con efecto en la evaluación.}
    \begin{itemize}
        \item El listado de ofertas usado para exportar el catálogo estaba limitado a 20 registros.
        \item La ingesta programada no podía desactivarse, lo que habría alterado el corpus congelado.
    \end{itemize}
\end{enumerate}

En términos de DSR, estos hallazgos son resultados de una evaluación \emph{formativa} del artefacto, realizada durante el ciclo de diseño y no después de él. Ilustran por qué la evaluación de un sistema existente debe comenzar por verificar que el artefacto implementa lo que su documentación describe.

\section{Iteración de diseño: recuperación}
\label{sec:art-iter-recuperacion}

\subsection{Diagnóstico}

Una primera corrida completa de la evaluación, con la configuración que resultó de la iteración anterior, y una revisión del pipeline frente a la literatura mostraron cuatro brechas entre lo que el trabajo mide y lo que el artefacto hacía:
\begin{itemize}
    \item \textbf{Las restricciones de la consulta no llegaban a la búsqueda vectorial.} Una consulta con señales de recuperación, como \textit{"find remote Python jobs paying over 150k"}, iba directamente a la búsqueda semántica y sus condiciones se ignoraban. Cuando el clasificador sí extraía las condiciones, la plataforma las resolvía con un filtro SQL que exigía la coincidencia de subcadena del título (\texttt{ILIKE}), es decir, la misma búsqueda léxica que el trabajo pretende superar, y pasaba al modelo hasta 30 ofertas ordenadas por fecha y no por relevancia. Además, el filtro no contemplaba la modalidad remota. En consecuencia, el benchmark de búsqueda filtrada (PE4) medía una capacidad que la aplicación no usaba.
    \item \textbf{El pipeline correspondía al esquema más básico de RAG.} En la clasificación de \citet{Gao2023}, el pipeline era un RAG \emph{ingenuo}: indexación, una única búsqueda densa y generación. Carecía de las optimizaciones posteriores a la recuperación que caracterizan al RAG \emph{avanzado}, como el reordenamiento, y de la recuperación híbrida que combina representaciones dispersas y densas.
    \item \textbf{Las consultas en español recibían respuestas en inglés.} En la primera corrida, 21 de las 25 consultas en español se respondieron en inglés: sin una instrucción explícita, el modelo respondía en el idioma del contexto, que es el de las ofertas.
    \item \textbf{El modelo de \textit{embeddings} era solo para inglés.} e5-base-v2 \citep{Wang2022E5} se entrenó con texto en inglés, de modo que las consultas en español dependían de la coincidencia accidental de términos.
\end{itemize}

\subsection{Cambios}

La Tabla~\ref{tab:art-cambios} resume los cambios de la iteración. El diseño de cada uno se describió en la Sección~\ref{sec:art-consulta}; aquí se resume su motivación.

\begin{table}[htbp!]
    \centering
    \caption{Cambios de la iteración de recuperación}
    \label{tab:art-cambios}
    \small
    \begin{tabular}{|p{3cm}|p{5.4cm}|p{5.4cm}|}
        \hline
        \textbf{Aspecto} & \textbf{Antes} & \textbf{Después} \\
        \hline
        Restricciones de la consulta & Ignoradas si la consulta tenía señales de recuperación; si no, filtro SQL por subcadena del título, con 30 ofertas ordenadas por fecha & Extraídas por el clasificador (inglés y español) y aplicadas como filtro dentro de la búsqueda densa y de BM25 \\
        \hline
        Primera etapa & Búsqueda densa de 15 fragmentos con umbral de similitud 0,5 & Búsqueda densa de 50 fragmentos y BM25 de 50 ofertas, fusionadas con RRF \\
        \hline
        Orden final & Similitud del mejor fragmento & Reordenamiento con Qwen3-Reranker-0.6B sobre 20 candidatas \\
        \hline
        Consultas sin ofertas relevantes & Siempre se devolvían cinco ofertas & Compuerta sobre el mejor puntaje del reordenamiento (0,134) \\
        \hline
        Modelo de \textit{embeddings} & e5-base-v2 (inglés) & multilingual-e5-base \\
        \hline
        Idioma de la respuesta & El del contexto & El de la pregunta \\
        \hline
        Salarios fuera de rango & Se almacenaban tal como llegaban & Se registran como no informados \\
        \hline
    \end{tabular}
\end{table}

\begin{description}
    \item[Filtros en la recuperación.] Aplicar el filtro durante la búsqueda, y no sobre los $k$ vecinos ya recuperados, evita que un filtro selectivo deje la lista incompleta, el problema que motiva los métodos de búsqueda aproximada con filtros \citep{Gollapudi2023, Patel2024}. Que la búsqueda densa y la léxica usen el mismo filtro garantiza que la fusión no reintroduzca ofertas que no cumplen las condiciones.
    \item[Búsqueda híbrida.] BM25 aporta la coincidencia exacta de términos poco frecuentes, como nombres de tecnologías, que un modelo denso puede diluir. Se eligió RRF \citep{Cormack2009} porque no requiere calibrar puntajes de naturaleza distinta ni datos de entrenamiento, con el valor habitual $k = 60$ y sin ajustarlo. \citet{Bruch2023} encontraron, sin embargo, que RRF es sensible a ese parámetro y que una combinación convexa de los puntajes la supera, con pocos ejemplos para ajustar su único peso.
    \item[Reordenamiento.] Un \textit{cross-encoder} lee la consulta y la oferta en una misma pasada, por lo que puede juzgar la relación entre ambas con más detalle que la comparación de dos vectores calculados por separado \citep{Nogueira2019}. Su costo crece con el número de candidatas, por eso solo se aplica a las 20 primeras de la lista fusionada.
    \item[Modelo de \textit{embeddings} multilingüe.] multilingual-e5-base \citep{Wang2024mE5} conserva la dimensión de 768 de e5-base-v2, de modo que el cambio no alteró el esquema. Los 4.013 fragmentos del corpus de evaluación se volvieron a codificar en el lugar, con los mismos identificadores y textos. Los vectores de e5-base-v2 se conservaron en una colección aparte para medir el efecto del cambio de modelo (Capítulo~\ref{cap:resultados}).
    \item[Idioma de la respuesta y salarios.] Los tres prompts de generación incluyen ahora la instrucción de responder en el idioma de la pregunta. La regla de salarios plausibles se aplicó también al corpus congelado: afectó a 12 ofertas, entre ellas un salario de 86.528.000 dólares para un puesto de despacho y montos en coronas suecas.
\end{description}

\subsection{Selección del modelo de reordenamiento}

La elección del modelo de reordenamiento, del número de candidatas y del umbral de la compuerta son decisiones de configuración. Tomarlas mirando las consultas de prueba sesgaría la evaluación a favor del artefacto, así que se tomaron con un conjunto de desarrollo separado y con una regla de decisión fijada antes de ver los resultados. En términos de DSR, es una evaluación formativa: sus resultados orientan el diseño y no se reportan como resultados del artefacto.

El conjunto de desarrollo contiene 20 necesidades de información, cinco por categoría, formuladas en 28 consultas en inglés y en español, sin solapamiento con las de prueba, más 10 consultas fuera de dominio que solo se usan para el umbral. Los juicios de relevancia siguieron el mismo procedimiento que la evaluación (Sección~\ref{sec:met-juicios}). Se compararon dos modelos de reordenamiento, bge-reranker-v2-m3, basado en BGE-M3 \citep{Chen2024M3}, y Qwen3-Reranker-0.6B \citep{Zhang2025Qwen3}, en cinco variantes: con y sin instrucción de tarea, y con el orden final dado solo por el reordenamiento o fusionado por RRF con el orden de la primera etapa. La regla tenía tres etapas:
\begin{enumerate}
    \item Con 30 candidatas, se elige la variante de mayor nDCG@10. Si otra queda a 0,02 o menos, se prefiere la más simple.
    \item Con la variante elegida, se usan 20 candidatas en lugar de 30 si el nDCG@10 baja 0,01 o menos.
    \item El umbral es el mayor valor que deja por debajo, como máximo, a una de las 28 consultas de desarrollo, evaluado sobre el mejor puntaje de cada consulta. Se adopta solo si deja por debajo al menos a 8 de las 10 consultas fuera de dominio.
\end{enumerate}

\begin{table}[htbp!]
    \centering
    \caption{Selección del modelo de reordenamiento en el conjunto de desarrollo (28 consultas, 30 candidatas)}
    \label{tab:art-rerank-dev}
    \begin{tabular}{|l|c|c|}
        \hline
        \textbf{Variante} & \textbf{nDCG@10} & \textbf{P@5} \\
        \hline
        Sin reordenamiento (densa + BM25 + RRF) & 0{,}742 & 0{,}757 \\
        \hline
        bge-reranker-v2-m3 & 0{,}804 & 0{,}814 \\
        \hline
        bge-reranker-v2-m3, fusionado con RRF & 0{,}798 & 0{,}800 \\
        \hline
        Qwen3-Reranker-0.6B & 0{,}820 & 0{,}864 \\
        \hline
        Qwen3-Reranker-0.6B con instrucción & 0{,}839 & 0{,}879 \\
        \hline
        Qwen3-Reranker-0.6B con instrucción, fusionado con RRF & 0{,}812 & 0{,}879 \\
        \hline
    \end{tabular}
\end{table}

La Tabla~\ref{tab:art-rerank-dev} muestra la primera etapa. Todas las variantes mejoraron entre 0,06 y 0,10 el nDCG@10 de la recuperación sin reordenamiento. Qwen3-Reranker-0.6B sin instrucción quedó a 0,019 de la mejor variante, por lo que la regla lo eligió por ser más simple. En la segunda etapa, pasar de 30 a 20 candidatas cambió el nDCG@10 de 0,832 a 0,830 y redujo la mediana de la latencia de recuperación de 1.184 a 788~ms, así que se adoptaron 20. En la tercera, el mejor puntaje más bajo entre las consultas de desarrollo fue 0,132 y el siguiente 0,134; las diez consultas fuera de dominio quedaron por debajo de 0,09. El umbral de 0,134 cumplió la condición y se adoptó.

Una observación cualitativa motivó incluir la fusión con RRF entre las variantes: bge-reranker-v2-m3 es muy literal. Asignó una probabilidad de 0,97 al par \textit{"security engineer"} y \textit{"Security Engineer, Cloud"}, pero de 0,002 al par formado por la descripción funcional \textit{"Cybersecurity roles defending company systems against attacks"} y una oferta de \textit{"Security Engineer"}. Se temía que perjudicara las consultas de desajuste de vocabulario; en el conjunto de desarrollo no ocurrió, y la variante elegida no usa la fusión.

\subsection{Defecto detectado durante la evaluación}

La corrida sobre el conjunto de prueba reveló un defecto de enrutamiento. Ante la consulta \textit{"what are the most recent job postings"}, el clasificador devolvía la etiqueta de recuperación junto con la plantilla de ofertas recientes, y la consulta se atendía por la ruta semántica. Se corrigió con la segunda regla de enrutamiento de la Sección~\ref{sec:art-clasificacion}: las plantillas de estadísticas, valores distintos y ofertas recientes se atienden siempre por la ruta de agregación. Las plantillas de listado conservan la etiqueta del clasificador, porque 51 de las 85 consultas de recuperación del conjunto de prueba llegan como recuperación acompañadas de una plantilla de listado por título o habilidad. El Capítulo~\ref{cap:resultados} reporta la exactitud de la ruta de agregación antes y después de la corrección.
